\documentclass[11pt,a4paper]{article}
\usepackage[T1]{fontenc}
\usepackage[utf8]{inputenc}
\usepackage{lmodern}
\usepackage{amsmath,amssymb,amsthm,mathtools,mathrsfs}
\usepackage{graphicx,float,xcolor,multicol}
\usepackage[shortlabels]{enumitem}
\usepackage[margin=27mm]{geometry}
\usepackage{microtype,needspace,placeins}
\usepackage{tikz,tikz-cd}
\usetikzlibrary{arrows.meta,calc,positioning,patterns,decorations.pathreplacing}
\usepackage[colorlinks=true,linkcolor=teal,urlcolor=teal,citecolor=teal]{hyperref}
\setlength{\parindent}{0pt}
\setlength{\parskip}{.6em}
\setcounter{secnumdepth}{0}
\allowdisplaybreaks
\raggedbottom
\providecommand{\cross}{\times}
\DeclareUnicodeCharacter{2020}{\ensuremath{\dagger}}
\DeclareMathOperator{\supp}{supp}
\DeclareMathOperator*{\esssup}{ess\,sup}
\providecommand{\chapter}{\section}
\newenvironment{tool}[2]{\subsection*{Tool #1 --- #2}}{}
\newcommand{\ToolRef}[1]{Tool #1}
\newcommand{\FigureTag}[2]{\textit{Figure #1}}
\newcommand{\Result}[1]{\subsection*{#1}}
\newcommand{\Application}[1]{\subsubsection*{Application\if\relax\detokenize{#1}\relax\else\space--- #1\fi}}
\newcommand{\ProofHeading}{\par\textit{Proof.}\space}
\newcommand{\ProofOf}[1]{\par\textit{Proof of #1.}\space}
\newcommand{\RemarkHeading}{\par\textbf{Remark.}\space}
\newcommand{\NoteHeading}{\par\textbf{Note.}\space}
\newcommand{\ExerciseHeading}{\par\textbf{Exercise.}\space}

% Rendering support only; mathematical text is in each independent post.
\usepackage{booktabs,array,longtable}
\definecolor{HandoutBlue}{HTML}{2F6690}
\definecolor{HandoutNavy}{HTML}{17365D}
\newtheorem{theorem}{Theorem}
\newtheorem{lemma}[theorem]{Lemma}
\newtheorem{proposition}[theorem]{Proposition}
\newtheorem{corollary}[theorem]{Corollary}
\newtheorem{conjecture}[theorem]{Conjecture}
\newtheorem{definition}{Definition}
\newtheorem{example}{Example}
\newtheorem{namedtoolcounter}{Tool}
\newenvironment{namedtool}[1]{\begin{namedtoolcounter}[#1]}{\end{namedtoolcounter}}
\newtheorem*{claim}{Claim}
\newtheorem*{remark}{Remark}
\newtheorem*{exercise}{Exercise}
\newtheorem*{question}{Question}
\newtheorem*{observation}{Observation}
\newcommand{\SourceChapter}[2]{\section*{#2}}
\newcommand{\Lecture}[2]{\section*{Lecture #1: #2}}
\newcommand{\unclear}[1]{\textit{[unclear: #1]}}
\newcommand{\sourcegap}{\textit{[unfinished in the source]}}
\DeclareMathOperator{\dist}{dist}
\DeclareMathOperator{\id}{id}
\DeclareMathOperator{\Hol}{Hol}
\DeclareMathOperator{\Per}{Per}
\DeclareMathOperator{\Fix}{Fix}
\DeclareMathOperator{\diam}{diam}
\DeclareMathOperator{\Var}{Var}
\DeclareMathOperator{\Leb}{Leb}
\DeclareMathOperator{\Hom}{Hom}
\DeclareMathOperator{\End}{End}
\DeclareMathOperator{\Ann}{Ann}
\DeclareMathOperator{\Spec}{Spec}
\DeclareMathOperator{\Max}{Max}
\DeclareMathOperator{\Ob}{Ob}
\DeclareMathOperator{\Tor}{Tor}
\DeclareMathOperator{\Ass}{Ass}
\DeclareMathOperator{\Mor}{Mor}
\DeclareMathOperator{\Fun}{Fun}
\DeclareMathOperator{\Nat}{Nat}
\DeclareMathOperator{\Cone}{Cone}
\DeclareMathOperator{\MCG}{MCG}
\DeclareMathOperator{\Aut}{Aut}
\DeclareMathOperator{\Deck}{Deck}
\DeclareMathOperator{\SL}{SL}
\DeclareMathOperator{\PSL}{PSL}
\DeclareMathOperator{\Area}{Area}
\DeclareMathOperator{\im}{im}
\DeclareMathOperator{\PSh}{PSh}
\newcommand{\CRing}{\mathsf{CRings}}
\newcommand{\AMod}{A\text{-}\mathsf{Mod}}
\newcommand{\Ab}{\mathsf{Ab}}
\newcommand{\Z}{\mathbb Z}
\newcommand{\Q}{\mathbb Q}
\newcommand{\R}{\mathbb R}
\newcommand{\C}{\mathbb C}
\newcommand{\N}{\mathbb N}
\newcommand{\kk}{\mathbb k}
\renewcommand{\aa}{\mathfrak a}
\newcommand{\bb}{\mathfrak b}
\newcommand{\pp}{\mathfrak p}
\newcommand{\qq}{\mathfrak q}
\newcommand{\mm}{\mathfrak m}
\newcommand{\nn}{\mathfrak n}
\newcommand{\rad}{\mathrm r}
\newcommand{\cat}[1]{\mathsf{#1}}
\setcounter{secnumdepth}{0}
\allowdisplaybreaks

\graphicspath{{figures/commutative-algebra/}}
\begin{document}
\section*{Primary Decomposition and Uniqueness}
% BLOG-CONTENT-BEGIN
\subsection{Primary decomposition}
\setcounter{definition}{8}\begin{definition}
An ideal $\qq<A$ is \emph{primary} if
\[
xy\in\qq\ \Longleftrightarrow\ x\in\qq\text{ or }y\in\qq\text{ or }x,y\in\rad(\qq).
\]
Equivalently, every zero-divisor in $A/\qq$ is nilpotent. If $\rad(\qq)=\pp$, then $\qq$ is \emph{$\pp$-primary}.
% Author check: the displayed equivalence is not silently replaced by the usual primary-ideal implication.
\end{definition}
\begin{claim}
$\rad(\qq)$ is the smallest prime ideal containing $\qq$.
\end{claim}
\begin{proof}
% Source page 6
If $xy\in\rad(\qq)$, then $x^ny^n\in\qq$ for $n>0$. Hence $x^n\in\qq$, $y^n\in\qq$, or $x^n,y^n\in\rad(\qq)$; thus $x$ or $y$ belongs to $\rad(\qq)$.
\end{proof}
\setcounter{example}{6}\begin{example}
\begin{enumerate}[label=(\roman*)]
\setcounter{enumi}{0}\item In $A=\Z$, $(0)$ and $(p^n)$ are primary ideals.
\setcounter{enumi}{1}\item In $A=\kk[x,y]$, let $\qq=(x,y^2)$. Then $A/\qq\cong\kk[y]/(y^2)$, whose zero-divisors are multiples of $y$. Hence $\qq$ is primary. Moreover, $\rad(\qq)=(x,y)=\pp$ is prime, and $\pp^2\subset\qq\subset\pp$. Thus not all primary ideals are powers of prime ideals.
\setcounter{enumi}{2}\item If $\rad(\aa)$ is maximal, then $\aa$ is primary. In particular, $\mm^n$ is $\mm$-primary. Here $A/\aa$ has only one prime ideal: $\rad(\aa)=\mm$ implies that $A/\aa\setminus\mm$ consists of units.
\end{enumerate}
\end{example}
\begin{remark}
In $\Z$, unique factorization gives $n=p_1^{a_1}\cdots p_n^{a_n}$, where $n\in\Z$ and the $p_i$ are primes. Equivalently,
\[(n)=\prod_{i=1}^n(p_i)^{a_i}=\bigcap_{i=1}^n(p_i)^{a_i}.\]
An appropriate generalization of this phenomenon to arbitrary rings requires replacing powers of primes by primary ideals. For certain rings, every ideal has an ``almost'' unique primary decomposition, $\aa=\bigcap_{i=1}^n\qq_i$, where $\qq_i$ is $\pp_i$-primary.
\end{remark}
% Source page 7
\setcounter{definition}{9}\begin{definition}
For an ideal $\aa<A$, an expression
\[\aa=\bigcap_{i=1}^n\qq_i,\qquad\qq_i\text{ is }\pp_i\text{-primary},\tag{*}\]
is a \emph{primary decomposition} of $\aa$. The set $\{\pp_i\}_{i=1}^n$ is the set of \emph{primes associated to $\aa$}.
\end{definition}
\begin{remark}
If $\qq_j\supseteq\bigcap_{k=1}^{\ell}\qq_{i_k}$, then $\qq_j$ is redundant in (*). Henceforth we assume our primary decomposition is \emph{irredundant}, or minimal.
\end{remark}
\setcounter{theorem}{8}\begin{lemma}
If the $\qq_i$ are $\pp$-primary, then so is $\bigcap_i\qq_i$.
\end{lemma}
\begin{proof}
We have $\rad(\qq)=\rad(\bigcap_i\qq_i)=\bigcap_i\rad(\qq_i)$.
% Source page 8
If $xy\in\qq$ and $y\notin\qq$, then there exists $i$ with $xy\in\qq_i$ and $y\notin\qq_i$, so $x\in\pp$.
% Author check: finiteness of the intersection is implicit in the surrounding decomposition.
\end{proof}
\paragraph{Consequence.}
If $\aa=\bigcap_{i=1}^n\qq_i$ is a primary decomposition, we can assume that the primes $\pp_i$ for which $\qq_i$ is $\pp_i$-primary are all distinct.
\setcounter{theorem}{9}\begin{lemma}
Let $\qq$ be $\pp$-primary and $x\in A$.
\begin{enumerate}[label=(\roman*)]
\setcounter{enumi}{0}\item If $x\in\qq$, then $(\qq:x)=A$.
\setcounter{enumi}{1}\item If $x\notin\qq$, then $(\qq:x)$ is $\pp$-primary, and consequently $\rad((\qq:x))=\pp$.
\end{enumerate}
\end{lemma}
\begin{proof}
(i) is clear. For (ii), $y\in(\qq:x)\Rightarrow xy\in\qq$. As $x\notin\qq$, we get $y\in\pp$. Thus
\[
\qq\subseteq(\qq:x)\subseteq\pp
\ \Longrightarrow\ \rad(\qq)\subseteq\rad((\qq:x))\subseteq\rad(\pp)
\ \Longrightarrow\ \rad((\qq:x))=\pp.
\]
If $yz\in(\qq:x)$ and $y\notin\pp$, then $xyz\in\qq$, so $xz\in\qq$ and hence $z\in(\qq:x)$.
\end{proof}
% Source page 9
\setcounter{theorem}{10}\begin{theorem}[First uniqueness theorem]
The associated primes in an irredundant primary decomposition depend only on $\aa$.
\end{theorem}
\begin{proof}
Note that
\[\rad((\aa:x))=\bigcap_i\rad((\qq_i:x))=\bigcap_{x\notin\qq_i}\pp_i.\]
If $\rad((\aa:x))$ is prime, then it equals $\pp_j$ for some $j$. Since the decomposition is irredundant, for every $1\le i\le n$ there exists $x_i\notin\qq_i$ with $x_i\in\bigcap_{j\ne i}\qq_j$. This yields $\rad((\aa:x_i))=\pp_i$. Therefore the $\pp_i$ are precisely the primes of the form $\rad((\aa:x))$, and hence depend on $\aa$.
\end{proof}
% Source page 10
\setcounter{theorem}{11}\begin{lemma}
Let $S$ be a multiplicatively closed subset of $A$ and $\qq$ a $\pp$-primary ideal.
\begin{enumerate}[label=(\roman*)]
\setcounter{enumi}{0}\item If $S\cap\pp\ne\varnothing$, then $S^{-1}\qq=S^{-1}A$.
\setcounter{enumi}{1}\item If $S\cap\pp=\varnothing$, then $S^{-1}\pp$-primary ideals of $S^{-1}A$ correspond to $\pp$-primary ideals of $A$.
\end{enumerate}
\end{lemma}
\begin{proof}
(i) If $s\in S\cap\pp$, then $s^n\in S\cap\qq$ for some $n>0$, so $S^{-1}\qq$ contains a unit.

(ii) If $\qq$ is $\pp$-primary, then $S^{-1}\qq$ is $S^{-1}\pp$-primary:
\[
\frac xs\frac yt\in S^{-1}\qq\ \Longrightarrow\ xy\in\qq
\ \Longrightarrow\ x\in\qq\text{ or }y^n\in\qq
\ \Longrightarrow\ \frac xs\in S^{-1}\qq\text{ or }\left(\frac yt\right)^n\in S^{-1}\qq.
\]
Also, $\rad(S^{-1}\qq)=S^{-1}\rad(\qq)=S^{-1}\pp$. Similarly, one can show that if $S^{-1}\qq$ is $S^{-1}\pp$-primary, then $\qq$ is $\pp$-primary.
% Author check: the final sentence has not been restricted to contracted ideals.
\end{proof}
% Source page 11
\subsection{Isolated and embedded primes}
\paragraph{Second uniqueness theorem.}
Let $\Sigma$ be a set of prime ideals associated to $\aa$. Say $\Sigma$ is \emph{isolated} if $\pp'\subseteq\pp\in\Sigma\Rightarrow\pp'\in\Sigma$. Set
\[S=A\setminus\bigcup_{\pp\in\Sigma}\pp.\]
Clearly $S$ is multiplicatively closed. Further, $\pp'\in\Sigma\Rightarrow\pp'\cap S=\varnothing$, while $\pp'\notin\Sigma\Rightarrow\pp'\nsubseteq\bigcup_{\pp\in\Sigma}\pp$ by prime avoidance, since $\pp'\nsubseteq\pp$ for every $\pp\in\Sigma$. Thus $\pp'\cap S\ne\varnothing$.

Let $\{\pp_i\}_{i=1}^n$ be the primes associated to $\aa$, and let $\{\pp_{i_1},\ldots,\pp_{i_k}\}$ be its minimal elements under inclusion. Then $\Sigma:=\{\pp_{i_j}\}_{j=1}^k$ is isolated. Hence
\[
(S^{-1}\aa)^c
=\bigcap_{i=1}^n(S^{-1}\qq_i)^c
=\bigcap_{j=1}^k(S^{-1}\qq_{i_j})^c
=\bigcap_{j=1}^k\qq_{i_j}.
\]
% Source page 12
Therefore $\bigcap_{j=1}^k\qq_{i_j}$ is determined only by $\aa$. The primes minimal among $\{\pp_i\}_{i=1}^n$ are called \emph{isolated primes}; the rest are \emph{embedded primes}.
\begin{remark}
One can generalize this to modules and prove analogous results. Briefly, for an $A$-module $M$, set
\[
\Ass(M):=\{\pp<A\text{ prime}\mid\pp=\Ann(m),\ m\in M\}.
\]
The primes associated to $\aa<A$ are elements of $\Ass(A/\aa)$. We say $N$ is a $\pp$-primary submodule of $M$ if $\Ass(M/N)=\{\pp\}$.
\end{remark}
Lasker proved, using methods from elimination theory, that affine rings admit primary decomposition for all ideals. Noether later developed the notion of chain conditions and extended Lasker's result to a wide range of rings satisfying these conditions. In her honour, these rings are named \emph{Noetherian rings}.
% BLOG-CONTENT-END
\end{document}
